% !TEX root = ../main.tex
\chapter{Resultados}

Este capítulo reúne la evidencia experimental de los experimentos de
reconstrucción, clasificación de WSI, segmentación de WSI a resolución de
parche y análisis geométrico del espacio latente. Se distingue entre el
conjunto fijo de validación usado para diagnóstico durante el desarrollo, las
pruebas selladas y los análisis exploratorios posteriores. Las diferencias se
expresan como VAE base menos VAE-\(\mathrm{SO}(2)\), salvo indicación contraria.

\section{Configuración experimental final}

Los resultados corresponden a los puntos de control de la actualización
\(60\,000\) de ambos VAE. Los codificadores y decodificadores permanecieron
congelados durante todas las tareas posteriores. La reconstrucción sellada
incluyó \(67\,138\) parches de 23 WSI; la clasificación de diagnóstico usó las
bolsas completas de esas mismas 23 WSI, con \(261\,168\) posiciones; y la
segmentación a resolución de parche evaluó \(31\,572\) posiciones de regiones
anotadas. Cada experimento supervisado tuvo una sola inicialización por rama,
por lo que los intervalos cuantifican remuestreo de WSI, no variación entre
entrenamientos.

La Figura~\ref{fig:entrenamiento-vae} resume las diez épocas de los dos VAE.
Ambas ramas completaron el presupuesto previsto. Las curvas describen el proceso de optimización.

\begin{figure}[p]
    \centering
    \includegraphics[width=0.88\textwidth]{vae_training_loss_panels.png}
    \caption{Evolución del objetivo compuesto y del término MAE durante
    \(60\,000\) actualizaciones. Las curvas de entrenamiento están suavizadas y
    los marcadores corresponden a las evaluaciones periódicas de validación.}
    \label{fig:entrenamiento-vae}
\end{figure}

\begin{figure}[p]
    \ContinuedFloat
    \centering
    \includegraphics[width=0.88\textwidth]{vae_training_quality_panels.png}
    \caption[]{Evolución del entrenamiento de los VAE (continuación): pérdida
    estructural \(1-\mathrm{SSIM}\) y PSNR sobre los 25 parches fijos de
    validación. }
\end{figure}

\begin{figure}[p]
    \ContinuedFloat
    \centering
    \includegraphics[width=0.88\textwidth]{vae_training_control_panels.png}
    \caption[]{Evolución del entrenamiento de los VAE (continuación): razón de
    equivarianza latente de desarrollo y calendario de la tasa de aprendizaje.
    La razón corresponde al seguimiento del desarrollo.}
\end{figure}

La Figura~\ref{fig:entrenamiento-vae-dispersion} complementa los diagnósticos
anteriores con una vista compacta de la pérdida de reconstrucción. La banda
muestra la desviación estándar registrada dentro de cada evaluación de
validación, no la variación entre entrenamientos independientes.

\begin{figure}[p]
    \centering
    \includegraphics[width=\textwidth]{vae_training_curves_refined.png}
    \caption{Pérdida de reconstrucción durante las 60.000 actualizaciones de
    cada VAE. Los puntos son medias de validación con corrupción determinista;
    la banda muestra la media más o menos una desviación estándar registrada.
    La línea punteada resume las pérdidas de entrenamiento por bloques de
    3.000 pasos. Cada arquitectura se entrenó una sola vez.}
    \label{fig:entrenamiento-vae-dispersion}
\end{figure}

\section{Reconstrucción y aprendizaje de representaciones}

\subsection{Inspección cualitativa y conjunto fijo de desarrollo}

Los dos modelos preservaron la organización general del tejido, los núcleos y
las variaciones principales de tinción en los 25 parches fijos. La
Figura~\ref{fig:reconstrucciones-fixed9} amplía los primeros nueve ejemplos
del selector y compara cada entrada con las reconstrucciones de ambos VAE.
La Figura~\ref{fig:reconstrucciones-fixed25}, situada a continuación,
conserva la vista general del conjunto fijo completo. Ambas figuras permiten inspeccionar los parches fijos de validación.

\begin{figure}[htbp]
    \centering
    \includegraphics[width=\textwidth]{vae_reconstructions_fixed9.png}
    \caption{Primeros nueve parches del selector fijo de validación, sus
    reconstrucciones con el VAE base y sus reconstrucciones con el
    VAE-\(\mathrm{SO}(2)\). Las tres cuadrículas de 3 por 3 conservan el mismo
    orden de parches y permiten observar más detalle que el mosaico completo.}
    \label{fig:reconstrucciones-fixed9}
\end{figure}

\begin{figure}[htbp]
    \centering
    \includegraphics[width=\textwidth]{vae_reconstructions_fixed25.png}
    \caption{Vista general de las entradas y reconstrucciones pareadas del
    conjunto fijo de 25 parches. Cada columna conserva el mismo parche para
    ambos VAE.}
    \label{fig:reconstrucciones-fixed25}
\end{figure}

La vista general permite revisar diversidad, pero comprime los detalles celulares.
La Figura~\ref{fig:reconstruccion-ampliada} muestra a resolución mayor el
parche fijo 12, que también se utilizó en los análisis de rotación. La
comparación ampliada ilustra, para este parche, que las diferencias entre modelos son sutiles en
la reconstrucción sin transformar y que ambas ramas suavizan parte de la
textura fina respecto de la entrada.

\begin{figure}[htbp]
    \centering
    \begin{subfigure}[t]{0.30\textwidth}
        \centering
        \includegraphics[width=\linewidth]{reconstruction_zoom_original.png}
        \caption{Parche original.}
    \end{subfigure}\hfill
    \begin{subfigure}[t]{0.30\textwidth}
        \centering
        \includegraphics[width=\linewidth]{reconstruction_zoom_normal.png}
        \caption{VAE base.}
    \end{subfigure}\hfill
    \begin{subfigure}[t]{0.30\textwidth}
        \centering
        \includegraphics[width=\linewidth]{reconstruction_zoom_so2.png}
        \caption{VAE-\(\mathrm{SO}(2)\).}
    \end{subfigure}
    \caption{Ampliación del parche fijo 12 y sus reconstrucciones finales sin
    transformar. Las tres imágenes se muestran a la misma escala y provienen
    de los arreglos numéricos conservados para el análisis de transformación
    latente.}
    \label{fig:reconstruccion-ampliada}
\end{figure}

En este conjunto fijo, el VAE base obtuvo MAE
\(0,0608\pm0,0127\), MSE \(0,0071\pm0,0028\), PSNR
\(27,898\pm1,839\) dB y SSIM \(0,7308\pm0,0637\). El
VAE-\(\mathrm{SO}(2)\) obtuvo, respectivamente,
\(0,0624\pm0,0130\), \(0,0075\pm0,0029\),
\(27,645\pm1,826\) dB y \(0,7170\pm0,0637\). Las dispersiones corresponden a
desviaciones estándar poblacionales entre los 25 parches.

\subsection{Prueba sellada de reconstrucción}

La Tabla~\ref{tab:reconstruccion-sellada} presenta las métricas sobre los
\(67\,138\) parches de prueba. En las cuatro medidas, el punto estimado favoreció
ligeramente al VAE base. La diferencia primaria de MAE fue \(-0,001358\), con
intervalo bootstrap por WSI del 95\% \([-0,002213;0,000046]\). Como el intervalo
incluyó cero, esta prueba no sustenta una diferencia concluyente en el criterio
primario. El VAE base tuvo menor MAE media en 22 de las 23 WSI.

\begin{table}[htbp]
    \centering
    \caption{Reconstrucción en la prueba sellada. Se reporta media
    \(\pm\) desviación estándar entre parches.}
    \label{tab:reconstruccion-sellada}
    \small
    \begin{tabular}{@{}lrrrr@{}}
        \toprule
        Modelo & MAE & MSE & PSNR (dB) & SSIM \\
        \midrule
        VAE base & \(0,0629\pm0,0239\) & \(0,0080\pm0,0074\) & \(27,923\pm2,721\) & \(0,7575\pm0,0842\) \\
        VAE-\(\mathrm{SO}(2)\) & \(0,0642\pm0,0223\) & \(0,0083\pm0,0067\) & \(27,658\pm2,595\) & \(0,7436\pm0,0862\) \\
        \bottomrule
    \end{tabular}
\end{table}

La Figura~\ref{fig:reconstruccion-sellada-wsi-distribuciones} resume las
medias de las 23 WSI para las cuatro métricas. Cada punto corresponde a una
lámina y las cajas resumen las distribuciones de medias por WSI. La
Figura~\ref{fig:reconstruccion-sellada-parches} conserva la distribución de
los parches y la Figura~\ref{fig:reconstruccion-sellada-wsi} presenta la
comparación pareada de MAE por WSI. PSNR y SSIM, definidos como medidas secundarias, también tuvieron mejores
valores descriptivos para el VAE base. En conjunto, se observó una diferencia
pequeña en la medida principal de esta prueba.

\begin{figure}[p]
    \centering
    \includegraphics[width=\textwidth]{vae_test_wsi_distributions.png}
    \caption{Distribuciones descriptivas de las medias de reconstrucción por
    WSI en las 23 láminas de prueba sellada. Cada punto es la media de los
    parches seleccionados de una WSI; las cajas resumen las distribuciones de las 23 medias por modelo.}
    \label{fig:reconstruccion-sellada-wsi-distribuciones}
\end{figure}

\begin{figure}[p]
    \centering
    \begin{subfigure}[t]{0.96\textwidth}
        \centering
        \includegraphics[width=\linewidth]{vae_test_patch_mae.png}
        \caption{Error absoluto medio en el dominio normalizado.}
    \end{subfigure}

    \begin{subfigure}[t]{0.96\textwidth}
        \centering
        \includegraphics[width=\linewidth]{vae_test_patch_mse.png}
        \caption{Error cuadrático medio en el dominio normalizado.}
    \end{subfigure}
    \caption{Distribuciones de error de reconstrucción por parche sobre las 23
    WSI de prueba. Las cajas resumen las distribuciones por parche; los parches de una
    misma lámina son observaciones correlacionadas.}
    \label{fig:reconstruccion-sellada-parches}
\end{figure}

\begin{figure}[p]
    \ContinuedFloat
    \centering
    \begin{subfigure}[t]{0.96\textwidth}
        \centering
        \includegraphics[width=\linewidth]{vae_test_patch_psnr.png}
        \caption{Relación señal a ruido pico por imagen.}
    \end{subfigure}

    \begin{subfigure}[t]{0.96\textwidth}
        \centering
        \includegraphics[width=\linewidth]{vae_test_patch_ssim.png}
        \caption{Índice de similitud estructural por imagen.}
    \end{subfigure}
    \caption[]{Distribuciones de calidad de reconstrucción por parche
    (continuación). Valores mayores indican mejor reconstrucción en PSNR y
    SSIM.}
\end{figure}

\begin{figure}[p]
    \centering
    \includegraphics[width=\textwidth]{vae_test_paired_wsi_mae_readable.png}
    \caption{MAE media pareada de las 23 WSI de prueba. Cada segmento une las medias de una misma WSI en el VAE base y el VAE-\(\mathrm{SO}(2)\).}
    \label{fig:reconstruccion-sellada-wsi}
\end{figure}

\section{Clasificación a nivel de WSI}

El clasificador MIL local-global se entrenó por separado, con la misma
arquitectura y protocolo, sobre cada representación congelada. La
Figura~\ref{fig:desarrollo-mil} muestra el desarrollo supervisado y el punto
seleccionado en cada rama. Las bolsas se procesaron completas, incluida la más
grande de \(32\,595\) parches.

\begin{figure}[p]
    \centering
    \includegraphics[width=\textwidth]{mil_training_f1.png}
    \caption{Evolución del entrenamiento y validación del clasificador MIL
    local-global: macro-F1 de validación. El punto de control de cada rama se
    eligió por esta métrica, con entropía cruzada como desempate. El azul
    identifica el VAE base y el naranja el VAE-\(\mathrm{SO}(2)\). }
    \label{fig:desarrollo-mil}
\end{figure}

\clearpage

\begin{figure}[p]
    \ContinuedFloat
    \centering
    \includegraphics[width=\textwidth]{mil_training_loss.png}
    \caption[]{Evolución del entrenamiento y validación del clasificador MIL
    local-global (continuación): entropía cruzada. Las líneas continuas
    corresponden a entrenamiento y las discontinuas a validación.}
\end{figure}

\clearpage

En la prueba sellada, el VAE-\(\mathrm{SO}(2)\) tuvo mejores estimaciones
puntuales de macro-F1, exactitud y entropía cruzada; el VAE base tuvo una
exactitud balanceada ligeramente mayor. La Tabla~\ref{tab:mil-sellado} muestra
que el intervalo de la diferencia primaria de macro-F1 atravesó cero. La
entropía cruzada fue menor para la rama equivariante y su intervalo exploratorio
no atravesó cero.

\begin{table}[htbp]
    \centering
    \caption{Clasificación de diagnóstico en 23 WSI de prueba sellada. La
    diferencia es VAE base menos VAE-\(\mathrm{SO}(2)\).}
    \label{tab:mil-sellado}
    \small
    \begin{tabular}{@{}lrrrr@{}}
        \toprule
        Métrica & VAE base & VAE-\(\mathrm{SO}(2)\) & Diferencia & IC 95\% \\
        \midrule
        Macro-F1 & 0,389 & 0,483 & \(-0,095\) & \([-0,207;0,050]\) \\
        Exactitud balanceada & 0,510 & 0,503 & \(0,007\) & \([-0,167;0,180]\) \\
        Exactitud & 0,522 & 0,609 & \(-0,087\) & \([-0,217;0,043]\) \\
        Entropía cruzada & 1,323 & 0,993 & \(0,330\) & \([0,135;0,524]\) \\
        \bottomrule
    \end{tabular}
\end{table}

La prueba contenía 5 WSI CC, 6 EC, 8 HGSC, 2 LGSC y 2 MC. La escasa presencia
de LGSC y MC hizo inestables sus medidas por clase. En particular, el VAE base
no recuperó casos EC ni LGSC, mientras que el VAE-\(\mathrm{SO}(2)\) no recuperó
casos MC. Las métricas globales se presentan en la
Figura~\ref{fig:mil-sellado-global}; las matrices de confusión, en la
Figura~\ref{fig:mil-sellado-confusiones}; y los F1 por clase, en la
Figura~\ref{fig:mil-sellado-f1-clases}.

\begin{figure}[p]
    \centering
    \includegraphics[width=\textwidth]{mil_test_metrics_refined.png}
    \caption{Macro-F1, exactitud balanceada y exactitud en las 23 WSI de
    prueba sellada. Cada par de barras presenta las puntuaciones de los dos modelos
    para una misma métrica; los intervalos por WSI se informan en la
    Tabla~\ref{tab:mil-sellado}.}
    \label{fig:mil-sellado-global}
\end{figure}

\begin{figure}[p]
    \centering
    \includegraphics[width=0.94\textwidth]{mil_test_confusion_normal.png}
    \caption{Matrices de confusión en la prueba sellada: clasificador entrenado
    sobre las representaciones del VAE base. Cada celda presenta el conteo y el
    porcentaje dentro de la clase real indicada por la fila. Ambos modelos se
    evaluaron sobre las mismas 23 WSI.}
    \label{fig:mil-sellado-confusiones}
\end{figure}

\begin{figure}[p]
    \ContinuedFloat
    \centering
    \includegraphics[width=0.94\textwidth]{mil_test_confusion_so2.png}
    \caption[]{Matrices de confusión en la prueba sellada (continuación):
    clasificador entrenado sobre las representaciones del
    VAE-\(\mathrm{SO}(2)\).}
\end{figure}

\clearpage

\begin{figure}[p]
    \centering
    \includegraphics[width=\textwidth]{mil_test_per_class_f1.png}
    \caption{F1 por subtipo histológico en la prueba sellada. El soporte
    pequeño de LGSC y MC limita la precisión de las estimaciones por clase.}
    \label{fig:mil-sellado-f1-clases}
\end{figure}

\clearpage

\section{Segmentación de WSI a resolución de parche}

La tarea de tejido asignó tumor, estroma o necrosis a cada celda seleccionada
de la cuadrícula de parches. Es, por tanto, una segmentación de WSI a resolución
de \(256\times256\) píxeles, pero no exhaustiva: solo se evaluaron regiones con
anotación de alta pureza. La prueba contenía \(21\,796\) parches de tumor,
\(8\,500\) de estroma y \(1\,276\) de necrosis. Esta última clase apareció en
solo cinco WSI.

La Tabla~\ref{tab:eficiencia-etiquetas} resume el macro-F1 para los cinco
presupuestos anidados. La rama equivariante tuvo una estimación mayor entre 250
y \(2\,500\) etiquetas por clase, mientras la rama base fue mayor con todas las
\(5\,671\) disponibles. Solo el intervalo simultáneo del presupuesto 500
excluyó cero. El área normalizada bajo la curva de eficiencia fue 0,615 para el
VAE base y 0,631 para el VAE-\(\mathrm{SO}(2)\), con diferencia
\(-0,016\;[-0,073;0,040]\); por consiguiente, la familia completa no demostró
una ventaja concluyente.

\begin{table}[htbp]
    \centering
    \caption{Macro-F1 en la segmentación dispersa de WSI a resolución de
    parche. Los intervalos simultáneos corresponden a la diferencia VAE base
    menos VAE-\(\mathrm{SO}(2)\).}
    \label{tab:eficiencia-etiquetas}
    \small
    \begin{tabular}{@{}rrrrr@{}}
        \toprule
        Etiquetas/clase & VAE base & VAE-\(\mathrm{SO}(2)\) & Diferencia & IC simultáneo 95\% \\
        \midrule
        250 & 0,508 & 0,523 & \(-0,015\) & \([-0,072;0,041]\) \\
        500 & 0,511 & 0,575 & \(-0,064\) & \([-0,121;-0,008]\) \\
        \(1\,000\) & 0,582 & 0,598 & \(-0,016\) & \([-0,072;0,041]\) \\
        \(2\,500\) & 0,701 & 0,714 & \(-0,013\) & \([-0,069;0,044]\) \\
        \(5\,671\) & 0,766 & 0,710 & \(0,055\) & \([-0,001;0,112]\) \\
        \bottomrule
    \end{tabular}
\end{table}

La Figura~\ref{fig:tejido-resultados} evidencia que la relación entre
representación y cantidad de etiquetas no fue monótona. También muestra la
dependencia por clase: necrosis tuvo el menor soporte y la mayor sensibilidad a
las cinco WSI que la contenían. Los resultados describen únicamente las
regiones anotadas y no asignan una clase a los sectores negros o sin máscara.

\clearpage

\begin{figure}[p]
    \centering
    \includegraphics[width=\textwidth]{tissue_training_f1.png}
    \caption{Mejor macro-F1 de validación alcanzado por el clasificador de
    tejido en cada presupuesto de etiquetas. El azul identifica el VAE base y
    el naranja el VAE-\(\mathrm{SO}(2)\). }
    \label{fig:tejido-desarrollo}
\end{figure}

\clearpage

\begin{figure}[p]
    \centering
    \includegraphics[width=\textwidth]{tissue_label_efficiency_refined.png}
    \caption{Macro-F1 de segmentación de WSI a resolución de parche en los
    cinco presupuestos balanceados y anidados de etiquetas por clase. El
    asterisco marca el presupuesto de 500 etiquetas cuyo intervalo simultáneo
    del 95\% para la diferencia entre modelos excluye cero.}
    \label{fig:tejido-resultados}
\end{figure}

\clearpage

\begin{figure}[p]
    \ContinuedFloat
    \centering
    \includegraphics[width=\textwidth]{tissue_test_per_class.png}
    \caption[]{Prueba sellada de segmentación a resolución de parche
    (continuación): desempeño de tumor, estroma y necrosis en los cinco
    presupuestos de etiquetas.}
\end{figure}

\clearpage

\section{Análisis del espacio latente}
\label{sec:resultados-latentes}

\subsection{Transformaciones prescritas y reconstrucciones}

La primera comprobación aplicó giros exactos de 90, 180 y 270 grados al mapa
latente espacial y decodificó el resultado. La Figura~\ref{fig:rotacion-latente}
muestra que el VAE-\(\mathrm{SO}(2)\) siguió la acción prescrita en estos cuartos
de vuelta, mientras que aplicar la misma operación geométrica al latente del
VAE base produjo cambios y artefactos no controlados.

\begin{figure}[htbp]
    \centering
    \includegraphics[width=\textwidth]{vae_rotation_0_90.png}
    \caption{Reconstrucciones al rotar la entrada y al transformar el latente.
    Parche fijo 12 a 0 y 90 grados; para cada ángulo se presentan el VAE base
    y el VAE-\(\mathrm{SO}(2)\). La primera columna muestra la entrada rotada;
    la segunda, su reconstrucción; la tercera, la decodificación del latente
    transformado. Los cuartos de vuelta aplican una reindexación exacta.
    Fuente: análisis de transformación latente.}
    \label{fig:rotacion-latente}
\end{figure}

\begin{figure}[p]
    \ContinuedFloat
    \centering
    \includegraphics[width=\textwidth]{vae_rotation_180_270.png}
    \caption[]{Reconstrucciones al rotar la entrada y al transformar el latente
    (continuación): giros exactos de 180 y 270 grados del mismo parche. Las filas
    y columnas siguen la disposición de la página anterior.}
\end{figure}

La ampliación de la Figura~\ref{fig:rotacion-latente-ampliada} aísla el giro de
90 grados y las dos intervenciones latentes. En el VAE base aparecen mezclas de
textura y contornos que no corresponden al objetivo rotado; en el
VAE-\(\mathrm{SO}(2)\), la transformación espacial del posterior conserva la
organización de la reconstrucción. Este ejemplo amplía el mecanismo resumido en el mosaico completo.

\begin{figure}[htbp]
    \centering
    \begin{subfigure}[t]{0.30\textwidth}
        \centering
        \includegraphics[width=\linewidth]{latent_action_zoom_target.png}
        \caption{Objetivo a 90 grados.}
    \end{subfigure}\hfill
    \begin{subfigure}[t]{0.30\textwidth}
        \centering
        \includegraphics[width=\linewidth]{latent_action_zoom_normal.png}
        \caption{Acción sobre el latente base.}
    \end{subfigure}\hfill
    \begin{subfigure}[t]{0.30\textwidth}
        \centering
        \includegraphics[width=\linewidth]{latent_action_zoom_so2.png}
        \caption{Acción sobre el latente \(\mathrm{SO}(2)\).}
    \end{subfigure}
    \caption{Ampliación del parche fijo 12 al transformar el latente 90 grados
    y decodificarlo. Las imágenes provienen de los arreglos numéricos del
    análisis de transformación latente.}
    \label{fig:rotacion-latente-ampliada}
\end{figure}

La evaluación cuantitativa distinguió entre transformar el latente y
canonizarlo antes de decodificar. Para rotaciones interpoladas de cinco grados,
la mediana de la razón de error de acción fue 0,713 en el VAE base y 0,550 en
el VAE-\(\mathrm{SO}(2)\); la razón de canonización fue 0,661 y 0,513,
respectivamente. La rama equivariante fue menor en los 25 parches. Para el control
exacto de cuartos de vuelta, la razón de acción descendió de 0,709453 en el VAE
base a \(2,18\times10^{-6}\) en el VAE-\(\mathrm{SO}(2)\), y la de
canonización de 0,710638 a 0,088854.

\begin{figure}[htbp]
    \centering
    \begin{subfigure}[t]{\textwidth}
        \centering
        \includegraphics[width=0.82\linewidth]{latent_population_ratios.png}
        \caption{Razones poblacionales para rotaciones interpoladas.}
    \end{subfigure}

    \begin{subfigure}[t]{\textwidth}
        \centering
        \includegraphics[width=0.96\linewidth]{latent_exact_d4.png}
        \caption{Control exacto a 90, 180 y 270 grados.}
    \end{subfigure}
    \caption{Respuesta de las reconstrucciones a transformaciones del latente.
    Valores menores indican mejor acuerdo con la acción o canonización
    prescrita. El control exacto evalúa la compatibilidad del decodificador equivariante con \(C_4\). Fuente: análisis de transformación latente.}
    \label{fig:razones-latentes}
\end{figure}

La distinción anterior resultó esencial. Al evaluar la composición completa
\(D(\boldsymbol{\mu}(R_\theta x))\) frente a
\(R_\theta D(\boldsymbol{\mu}(x))\), que también incluye al codificador y la
interpolación de la entrada, la mediana fue 0,174 para el VAE base y 0,210 para
el VAE-\(\mathrm{SO}(2)\); el VAE base fue menor en los 25 parches.

\subsection{Geometría local y organización espacial}

En órbitas muestreadas cada grado, la medida corregida de geometría local fue
0,649 para el VAE base y 0,601 para el VAE-\(\mathrm{SO}(2)\), una reducción
relativa de 7,38\%. Los 25 parches favorecieron a la rama equivariante a esa
resolución. La dirección cambió al usar pasos de cinco grados: el VAE-\(\mathrm{SO}(2)\)
fue favorable en solo 8 de 25 parches.

La Figura~\ref{fig:orbita-latente-parche} muestra cómo se obtuvo esa evidencia
para el parche fijo 12. Cada punto representa el posterior completo dentro del
disco central para un ángulo, proyectado sobre un PCA ajustado con las 360
vistas del mismo modelo y parche. La trayectoria del VAE-\(\mathrm{SO}(2)\) es
visualmente más regular en este ejemplo; la curva inferior cuantifica el residual directo para cada ángulo.

\begin{figure}[htbp]
    \centering
    \includegraphics[width=0.96\textwidth]{latent_orbit_single_patch.png}
    \caption{Órbita latente de un parche al rotar la entrada entre 0 y 359
    grados. Arriba se muestran los dos PCA independientes; abajo, el residual
    directo para el parche y el control exacto agregado de cuartos de vuelta.
    Fuente: análisis de órbitas latentes.}
    \label{fig:orbita-latente-parche}
\end{figure}

La proyección PCA espacial de la Figura~\ref{fig:pca-latente-espacial} mostró
estructura local en ambos modelos. La RMS relativa de aristas tuvo mediana
1,4567 para el VAE base y 1,4676 para el VAE-\(\mathrm{SO}(2)\); por tanto, esta
medida no mostró mayor suavidad espacial de la rama equivariante. Un sondeo RGB
puntual explicó \(R^2=0,286\) y \(R^2=0,189\), respectivamente.

\begin{figure}[htbp]
    \centering
    \includegraphics[width=0.70\textwidth]{latent_spatial_pca_maps.png}
    \caption{Proyección PCA de los mapas latentes espaciales para cuatro
    parches fijos. Cada color representa los tres primeros componentes bajo un
    ajuste común por modelo; se usa para visualizar la organización interna. Fuente: análisis de órbitas y
    estructura latente.}
    \label{fig:pca-latente-espacial}
\end{figure}

\begin{figure}[htbp]
    \ContinuedFloat
    \centering
    \begin{subfigure}[t]{\textwidth}
        \centering
        \includegraphics[width=0.96\linewidth]{latent_spatial_pca_paired.png}
        \caption{RMS relativa de aristas, pareada sobre los 25 parches.}
    \end{subfigure}

    \begin{subfigure}[t]{\textwidth}
        \centering
        \includegraphics[width=0.96\linewidth]{latent_spatial_pca_delta.png}
        \caption{Diferencia VAE-\(\mathrm{SO}(2)\) menos VAE base por parche.}
    \end{subfigure}
    \caption[]{Proyección PCA espacial (continuación): medidas cuantitativas de
    rugosidad local. Valores menores indican menor rugosidad local.}
\end{figure}

El análisis latente verificó la compatibilidad numérica del decodificador
VAE-\(\mathrm{SO}(2)\) con la acción discreta \(C_4\) en cuartos de vuelta
y mostró una mejora local descriptiva a un grado. La comparación continua de
extremo a extremo favoreció al VAE base.

\clearpage

\section{Síntesis de resultados por objetivo}

La Tabla~\ref{tab:sintesis-objetivos} relaciona los objetivos específicos con
la evidencia disponible, indicando los resultados principales y el alcance de
cada evaluación.

\begin{table}[htbp]
    \centering
    \caption{Síntesis de evidencia por objetivo específico.}
    \label{tab:sintesis-objetivos}
    \small
    \begin{tabularx}{\textwidth}{@{}p{0.20\textwidth}XX@{}}
        \toprule
        Objetivo & Evidencia y resultado principal & Limitación principal \\
        \midrule
        Selección de datos & UBC-OCEAN permitió entrenamiento no supervisado,
        diagnóstico de cinco subtipos y evaluación de tres tejidos. & Las
        máscaras son suplementarias y no cubren exhaustivamente las WSI. \\
        Formulación e implementación & Se construyeron VAE residuales pareados,
        uno con convoluciones restringidas a \(F_0/F_1\), y clasificadores
        idénticos sobre sus latentes congelados. & La rama equivariante tiene
        menos parámetros y mayor costo de convolución. \\
        Evaluación & La reconstrucción favoreció levemente al VAE base; WSI y
        algunos presupuestos bajos favorecieron puntualmente al
        VAE-\(\mathrm{SO}(2)\); se verificó su control exacto \(C_4\). & Una
        inicialización por rama, 23 WSI de prueba, clases minoritarias escasas y
        análisis latente sobre 25 parches fijos. \\
        \bottomrule
    \end{tabularx}
\end{table}
